% ============================================================================
%  Why auxiliary-calibrated label-recovery attacks collapse on well-trained
%  models, and how a scale-invariant estimator avoids it.
%
%  Compile:  pdflatex VLI_report.tex   (run twice for cross-references)
%  Packages used are all part of a standard TeX Live / MiKTeX installation.
% ============================================================================
\documentclass[11pt]{article}

\usepackage[margin=1in]{geometry}
\usepackage[T1]{fontenc}
\usepackage{lmodern}
\usepackage{amsmath,amssymb,amsthm,mathtools}
\usepackage{booktabs}
\usepackage[protrusion=true,expansion=false]{microtype}
\usepackage[hidelinks]{hyperref}

\newtheorem{assumption}{Assumption}
\newtheorem{lemma}{Lemma}
\newtheorem{theorem}{Theorem}
\newtheorem{proposition}{Proposition}
\newtheorem{corollary}{Corollary}
\theoremstyle{remark}
\newtheorem{remark}{Remark}

\newcommand{\R}{\mathbb{R}}
\newcommand{\E}{\mathbb{E}}
\newcommand{\one}{\mathbf{1}}
\newcommand{\pos}[1]{\left[#1\right]_+}
\newcommand{\simplex}{\Delta_n}
\DeclareMathOperator*{\argmin}{arg\,min}
\DeclareMathOperator*{\argmax}{arg\,max}
\DeclareMathOperator{\softmax}{softmax}
\DeclareMathOperator{\LSE}{LSE}

\title{Label Recovery from Local Updates:\\
Why Auxiliary-Calibrated Attacks Collapse on Well-Trained Models,\\
and a Scale-Invariant Alternative}
\author{}
\date{\today}

\begin{document}
\maketitle

\begin{abstract}
We study the recovery of per-class label counts in one client window of
federated learning from the accumulated cross-entropy gradient of the final
linear layer. We prove two exact identities relating that gradient to
per-class error mass, and use them to isolate a single failure mechanism
shared by the published baselines iLRG, LLG$^+$, ZLG$^+$ and RLU: each
combines the observed update with a statistic of \emph{absolute magnitude}
measured outside the attacked window. Since the gradient norm is bounded by
the training error (Lemma~\ref{lem:norm}) while the held-out error stays
bounded away from zero (Proposition~\ref{prop:gap}), the two magnitudes
diverge as the model improves, and the estimator converges to a quantity
that does not depend on the attacked labels at all
(Theorem~\ref{thm:collapse}). We then present VLI and its closed form
VLI$^*$, whose defining property is invariance to the scale of the update
(Proposition~\ref{prop:invariance}); under a symmetry assumption that we show
is the unique label-exchangeable choice (Proposition~\ref{prop:exchange}),
VLI$^*$ reduces to a single scale-free expression that is exact
(Theorem~\ref{thm:closed}). We state the assumption's failure modes
quantitatively rather than claiming it holds.
\end{abstract}

\tableofcontents

% ============================================================================
\section{Scope and contribution}

\paragraph{What this report establishes.}
\begin{enumerate}
\item Two exact identities for the final-layer bias gradient
  (Lemmas~\ref{lem:norm} and~\ref{lem:identities}), including the
  zero-sum property $\sum_j G_j = 0$ that later does most of the work.
\item A single theorem (Theorem~\ref{thm:collapse}) that explains, in one
  mechanism, the degradation of every auxiliary-calibrated baseline as the
  global model improves, together with three supporting propositions
  (held-out error gap, Jensen gap in logit averaging, non-Lipschitz
  renormalization).
\item A scale-invariance property (Proposition~\ref{prop:invariance}) that
  the proposed estimator satisfies and the baselines do not, and a closed
  form (Theorem~\ref{thm:closed}) that is exact under a stated symmetry.
\item A uniqueness argument (Proposition~\ref{prop:exchange}) replacing the
  symmetry assumption's usual justification ``for simplicity'' with
  ``the unique choice available to an attacker with no class-specific
  knowledge''.
\item A sufficient condition strictly weaker than that symmetry
  (Proposition~\ref{prop:weak}), which widens the region of exactness.
\end{enumerate}

\paragraph{What this report does \emph{not} establish.}
The difficulty vector $r$ and confusion matrix $Q$ introduced below are
\emph{not} claimed to be recoverable from a shared update; the generalized
estimator of Theorem~\ref{thm:general} takes them as inputs and is used here
as an oracle upper bound. The numerical tables in
Section~\ref{sec:numerics} are evaluations on the generative model implied by
the assumptions, not experiments on a trained network; they verify algebra
and probe sensitivity, and are labelled as such throughout.

% ============================================================================
\section{Notation}
\label{sec:notation}

A client window contains $n$ examples drawn from $C$ classes. For
example $i$:

\begin{center}
\begin{tabular}{ll}
\toprule
Symbol & Meaning \\
\midrule
$y_i \in \{1,\dots,C\}$ & true label of example $i$ \\
$z_i \in \R^C$ & logits, $z_i = W e_i + b$ where $e_i$ is the embedding \\
$p_i = \softmax(z_i)$ & predicted class distribution \\
$e_{y_i}$ & one-hot indicator of the true class \\
$\ell_i = -\log p_{i,y_i}$ & cross-entropy loss \\
$L_{\mathrm{train}} = n^{-1}\sum_i \ell_i$ & mean training loss on the window \\
\midrule
$N_j = \#\{i : y_i = j\}$ & true count of class $j$; $\sum_j N_j = n$ \\
$\pi_j = N_j/n$ & true class proportion \\
$P_j = \sum_i p_{i,j}$ & total predicted mass on class $j$ \\
$G = \sum_{i=1}^n (p_i - e_{y_i})$ & \textbf{bias-gradient signal} \\
\midrule
$\epsilon_i = 1 - p_{i,y_i}$ & per-example error \\
$\bar\epsilon = n^{-1}\sum_i \epsilon_i$ & mean error on the window \\
$\bar\epsilon_j = N_j^{-1}\sum_{i:y_i=j}\epsilon_i$ & mean error of class $j$ (present $j$) \\
$a_j = N_j \bar\epsilon_j$ & \textbf{error mass} of class $j$; $a_j = 0$ if absent \\
$r_j = \bar\epsilon_j/\bar\epsilon$ & relative difficulty; $a_j = N_j r_j \bar\epsilon$ \\
$\kappa = r_{\max}/r_{\min}$ & difficulty spread over present classes \\
\midrule
$E_j = \sum_{i:y_i\ne j} p_{i,j}$ & error mass flowing \emph{into} class $j$ \\
$Q_{jk}$ & fraction of class $j$'s error mass sent to $k$; $Q_{jj}=0$ \\
$v$ & Perron vector of $Q^T$, $Q^Tv = v$, $\one^Tv=1$ \\
\midrule
$M_i = z_{i,y_i} - \LSE_{k\ne y_i} z_{i,k}$ & margin; $\epsilon_i = (1+e^{M_i})^{-1}$ \\
$\simplex = \{x \in \R^C : x \ge 0,\ \one^Tx = n\}$ & scaled probability simplex \\
$\pos{x} = \max(x,0)$ & componentwise positive part \\
\bottomrule
\end{tabular}
\end{center}

The off-diagonal error-transition matrix is defined for $a_j > 0$ by
\[
 Q_{jk} = \frac{\sum_{i : y_i = j} p_{i,k}}{a_j} \quad (k \ne j),
 \qquad Q_{jj} = 0 .
\]
Every defined row is nonnegative and sums to one, since
$\sum_{k\ne j}\sum_{i:y_i=j} p_{i,k} = \sum_{i:y_i=j}(1-p_{i,j}) = a_j$.
A class with zero error mass may be assigned any distribution over the other
classes; it contributes nothing to $Q^Ta$. Dividing by $a_j$ is what removes
$\bar\epsilon$ from $Q$: the matrix carries the \emph{shape} of the model's
confusion, not its magnitude. The same holds for $r$, a ratio of errors.
This distinction between \emph{shape} and \emph{magnitude} is the pivot of
the entire report.

\paragraph{Signal convention.}
An implementation accumulates per-minibatch \emph{mean}-loss gradients. With
all minibatches of size $B$ and $m$ steps per window ($n = mB$), the
accumulated signal differs from the example-sum $G$ above by a common
positive factor, and possibly by a further common window normalization. Every
formula below that is declared scale-invariant is unaffected by such a
factor. Under plain SGD a server weight update has the opposite sign; $G$
always denotes the loss gradient.

% ============================================================================
\section{Exact identities}

\begin{lemma}[Gradient norm is controlled by training error]
\label{lem:norm}
For every example, $\|p_i - e_{y_i}\|_1 = 2\epsilon_i$. Consequently
\[
 \|G\|_1 \;\le\; 2n\bar\epsilon \;\le\; 2n L_{\mathrm{train}} .
\]
\end{lemma}

\begin{proof}
The true-class coordinate of $p_i - e_{y_i}$ has absolute value
$1 - p_{i,y_i} = \epsilon_i$. The remaining coordinates are nonnegative and
sum to $1 - p_{i,y_i} = \epsilon_i$. Adding gives $2\epsilon_i$. The
triangle inequality yields
$\|G\|_1 \le \sum_i \|p_i - e_{y_i}\|_1 = 2n\bar\epsilon$. Finally
$1 - x \le -\log x$ for $x \in (0,1]$; apply it at $x = p_{i,y_i}$ and
average.
\end{proof}

Lemma~\ref{lem:norm} is the quantitative statement of the phenomenon this
report explains: \emph{as the model fits its training data, the entire signal
available to any attacker shrinks proportionally to the training error.} It
does not say the label information disappears --- Lemma~\ref{lem:identities}
shows the information survives in the \emph{direction} of $G$ --- but it does
say that any procedure comparing $G$ against a fixed external magnitude will
eventually be dominated by that external quantity.

\begin{lemma}[Count, error-mass and confusion identities]
\label{lem:identities}
Let $P_j = \sum_i p_{i,j}$. Then
\[
 G_j = P_j - N_j = -a_j + E_j,
 \qquad N_j = P_j - G_j,
 \qquad E_j = (Q^Ta)_j ,
\]
and therefore
\begin{equation}
 G = -(I - Q^T)\,a .
 \label{eq:system}
\end{equation}
Moreover
\begin{equation}
 \sum_{j=1}^C G_j = 0 ,
 \label{eq:zerosum}
\end{equation}
and if $N_j = 0$ then $G_j = P_j \ge 0$.
\end{lemma}

\begin{proof}
The $j$th coordinate of $G$ accumulates all predicted mass on class $j$ and
subtracts one per true class-$j$ example, giving $P_j - N_j$. Split $P_j$
into its true-class part $\sum_{i:y_i=j}p_{i,j} = N_j - a_j$ and its
other-class part $E_j$; this yields $G_j = -a_j + E_j$ and the absent-class
sign statement. For the confusion form, for each target $k$,
\[
 E_k = \sum_{j \ne k}\ \sum_{i : y_i = j} p_{i,k}
     = \sum_j a_j Q_{jk} = (Q^Ta)_k ,
\]
and substituting into $G = -a + E$ gives~\eqref{eq:system}. For
\eqref{eq:zerosum}, $\sum_j P_j = \sum_i \sum_j p_{i,j} = n$ and
$\sum_j N_j = n$.
\end{proof}

\begin{remark}
Identity~\eqref{eq:zerosum} has three consequences used repeatedly below.
(i) A nonzero $G$ must have a strictly positive coordinate, so the anchor
used by VLI$^*$ exists whenever the signal is not identically zero.
(ii) Mass conservation in Theorem~\ref{thm:closed} is automatic, so no
renormalization step is required. (iii) Each minibatch mean gradient is
itself zero-sum, since $B^{-1}(B - B) = 0$; the accumulated signal therefore
inherits~\eqref{eq:zerosum} up to floating-point error, even though the model
moves between steps.
\end{remark}

\begin{remark}[Validity across local steps]
Lemmas~\ref{lem:norm} and~\ref{lem:identities} are per-example statements and
hold when different examples are evaluated at different local model states.
They do \emph{not} automatically describe a server's optimizer update when
that update also contains momentum state, weight decay, a proximal term, or a
FedDyn correction. Section~\ref{sec:limits} returns to this.
\end{remark}

% ============================================================================
\section{Baseline estimators}
\label{sec:baselines}

We write each published baseline in the form in which it is implemented, so
that the next section can identify precisely which quantity each one borrows
from outside the attacked window. Throughout, $g = G/n$ denotes the
normalized signal and quantities carrying a bar or a hat
($\bar p$, $\bar O$, $\mathrm{shift}$, $\mathrm{offset}$) are measured on an
\emph{auxiliary} dataset available to the attacker --- in practice a
stratified fraction of the test set, hence approximately class-balanced.

\paragraph{iLRG (Ma et al., 2023).}
Recovers a class-wise embedding estimate as the ratio of the weight and bias
gradients, $\hat e_j = \nabla_{W_j}\!/\,\nabla_{b_j}$, feeds it through the
classifier to obtain recovered probability vectors, and solves a linear
system for the counts by pseudo-inverse.

\paragraph{LLG (Wainakh et al., 2021).}
Computes a single scalar \emph{impact} from the observed signal itself,
\[
 \mathrm{impact} = \Bigl(1 + \tfrac1C\Bigr)\frac{\sum_{j : G_j<0} G_j}{B},
\]
selects every class with a negative coordinate, and then greedily assigns the
remaining slots to the current $\argmin$ coordinate, decrementing by
$\mathrm{impact}$ after each assignment.

\paragraph{LLG$^+$ (Wainakh et al., 2021).}
As LLG, but first subtracts an auxiliary-measured offset,
\[
 G \longleftarrow G - \mathrm{offset}, \qquad
 \mathrm{offset}_j = \text{mean non-target coordinate $j$ on auxiliary data},
\]
and averages the auxiliary impact with the one measured before training.

\paragraph{ZLG$^+$ (Geng et al., 2021).}
\begin{equation}
 \widehat N_j = n\Bigl(\bar p_j - \frac{G_j}{\bar O}\Bigr),
 \label{eq:zlg}
\end{equation}
followed by clipping at zero and a multiplicative rescaling to total $n$;
$\bar p$ is an auxiliary posterior and $\bar O$ an auxiliary mean embedding
norm. \emph{Provenance note:} estimating $\bar p$ and $\bar O$ from a real
auxiliary set is not Geng et al.'s own design --- their paper estimates
both quantities by passing \emph{dummy}, synthetically generated images
through the model (a zero-shot construction, consistent with their stated
motivation of requiring no external dataset). The real-auxiliary-data form
used here and throughout this report follows the ``ZLG$^+$'' baseline
introduced by Chen \& Vikalo (RLU) for a fair comparison across methods,
all drawing on the same auxiliary set $A$; the dummy-data original is
denoted plain ZLG where the two are distinguished.

\paragraph{RLU (Chen and Vikalo, 2024).}
\begin{equation}
 \widehat N = n \cdot \frac{\pos{\mathrm{shift}\cdot\bar\rho - G}}
                           {\bigl\|\pos{\mathrm{shift}\cdot\bar\rho - G}\bigr\|_1},
 \qquad \mathrm{shift} = \softmax(\mu_{\mathrm{aux}}),
 \label{eq:rlu}
\end{equation}
where $\mu_{\mathrm{aux}}$ is the vector of \emph{mean logits} per class
measured on auxiliary data, $\bar\rho$ is a scalar accumulated-step
coefficient, and a posterior search then perturbs $\widehat N$ in units of
$m$ to better match $\mu_{\mathrm{aux}}$.

\begin{center}
\begin{tabular}{llll}
\toprule
Method & External quantity used & Its units & Form \\
\midrule
iLRG      & none (uses $\nabla W/\nabla b$) & --- & pseudo-inverse \\
LLG       & none                            & --- & self-calibrated greedy \\
LLG$^+$   & $\mathrm{offset}$               & absolute, $\propto \bar\epsilon_{\mathrm{aux}}$ & additive \\
ZLG$^+$   & $\bar p$, $\bar O$              & absolute                                        & \eqref{eq:zlg} \\
RLU       & $\mathrm{shift}$                & absolute, $\propto \bar\epsilon_{\mathrm{aux}}$ & \eqref{eq:rlu} \\
\bottomrule
\end{tabular}
\end{center}

% ============================================================================
\section{The failure mechanism}
\label{sec:failure}

\subsection{Auxiliary-only calibration becomes label-independent}

\begin{theorem}[Collapse of auxiliary-calibrated estimators]
\label{thm:collapse}
Let an estimator of normalized class proportions take the form
\[
 \widehat\pi = \Phi\bigl(c - G/n\bigr),
\]
where $c$ depends only on information outside the attacked client window and
$\Phi$ is $L_\Phi$-Lipschitz in a stated norm. Then
\begin{equation}
 \bigl\|\widehat\pi - \Phi(c)\bigr\|
 \;\le\; \frac{L_\Phi}{n}\|G\|_1
 \;\le\; 2L_\Phi\,\bar\epsilon
 \;\le\; 2L_\Phi\,L_{\mathrm{train}} .
 \label{eq:collapse}
\end{equation}
Hence as $L_{\mathrm{train}} \to 0$ the estimate converges to $\Phi(c)$, a
vector determined entirely by $c$ and not by the attacked labels. For an
estimator returning counts $n\Phi(\cdot)$ the bound is multiplied by $n$.
\end{theorem}

\begin{proof}
Lipschitz continuity gives
$\|\Phi(c - G/n) - \Phi(c)\| \le L_\Phi \|G/n\|_1$; apply
Lemma~\ref{lem:norm}.
\end{proof}

Both~\eqref{eq:zlg} and~\eqref{eq:rlu} are instances of this form, with
$c = \bar p$ and $c = \mathrm{shift}\cdot\bar\rho$ respectively, and
LLG$^+$'s offset subtraction is the same construction applied before a greedy
rule. The limit $\Phi(c)$ is informative:

\begin{itemize}
\item For ZLG$^+$, $\Phi(c) = n\bar p$. Since the auxiliary set is built
  class-balanced, $\bar p$ is approximately uniform, so the estimator
  degenerates to \emph{predicting a uniform class distribution}. Under a
  Dirichlet partition the overlap between the true client distribution and
  the uniform one is approximately $0.54$ at $\alpha = 0.5$ and $0.29$ at
  $\alpha = 0.1$ for $C = 10$; these are the instance-recovery values the
  method approaches. The failure is predictable, not random.
\item For RLU, $\Phi(c) = n\,\mathrm{shift}/\|\mathrm{shift}\|_1$, again a
  label-independent vector.
\item For LLG$^+$, the signal $G$ is overwhelmed by $-\mathrm{offset}$ and
  the subsequent $\argmin$ ordering is governed by the offset, not by the
  window. This predicts the counter-intuitive conclusion that LLG$^+$ is
  \emph{worse} than LLG, whose impact is computed from $G$ itself.
\end{itemize}

\begin{remark}[Limitation of the theorem]
Theorem~\ref{thm:collapse} applies to an implementation only after its exact
map, clipping and rounding have been shown to satisfy the representation and
the Lipschitz condition. It is not by itself a proof that every branch of
RLU's posterior search satisfies~\eqref{eq:collapse}. Indeed
Proposition~\ref{prop:lip} shows that RLU's normalization is \emph{not}
globally Lipschitz, so for that method the theorem should be read as
identifying the mechanism rather than as a tight bound.
\end{remark}

\subsection{Why the external magnitude does not shrink with the signal}

\begin{proposition}[Held-out error gap]
\label{prop:gap}
If an example is misclassified by the $\argmax$ rule, then
$p_{i,y_i} \le 1/2$ and hence $\epsilon_i \ge 1/2$. Consequently, if the
auxiliary error rate is $\tau_{\mathrm{err}}$,
\[
 \bar\epsilon_{\mathrm{aux}} \;\ge\; \tau_{\mathrm{err}}/2 .
\]
If $\bar\epsilon_{\mathrm{train}} \to 0$ while $\tau_{\mathrm{err}}$ stays
bounded away from zero, then
$\bar\epsilon_{\mathrm{aux}}/\bar\epsilon_{\mathrm{train}} \to \infty$.
\end{proposition}

\begin{proof}
On a misclassified example some $k \ne y_i$ has $p_{i,k} \ge p_{i,y_i}$;
since $p_{i,k} + p_{i,y_i} \le 1$ we get $p_{i,y_i} \le 1/2$. Average the
resulting lower bound on $1 - p_{i,y_i}$ over the auxiliary set, discarding
the nonnegative contribution of correctly classified examples.
\end{proof}

Proposition~\ref{prop:gap} is the second half of the mechanism. Lemma
\ref{lem:norm} says $\|G\| \to 0$; Proposition~\ref{prop:gap} says the
auxiliary magnitudes entering $c$ do \emph{not}. Their ratio diverges, so the
subtraction $c - G/n$ in Theorem~\ref{thm:collapse} is a comparison between
two quantities on incommensurable scales. \emph{Generalization gap, not
estimator design, is what breaks these attacks.}

\subsection{Two further defects specific to RLU}

\begin{proposition}[Jensen gap in logit averaging]
\label{prop:jensen}
For a sample with margin $M$, $\epsilon = (1+e^M)^{-1}$, and
$\epsilon \approx e^{-M}$ when $M$ is large and positive. If within-class
margins are approximately $\mathcal N(\mu,\sigma^2)$ and the large-margin
approximation is valid over most of their mass, then
\[
 \E\bigl[e^{-M}\bigr] = e^{-\mu + \sigma^2/2},
 \qquad e^{-\E[M]} = e^{-\mu},
\]
so substituting a mean margin understates this surrogate by the factor
$e^{\sigma^2/2}$. Jensen's inequality gives
$\E[e^{-M}] \ge e^{-\E[M]}$ unconditionally.
\end{proposition}

Equation~\eqref{eq:rlu} sets $\mathrm{shift} = \softmax(\mu_{\mathrm{aux}})$,
a softmax of \emph{mean logits}, i.e.\ exactly the substitution
Proposition~\ref{prop:jensen} describes. The bias is systematic and in the
direction of \emph{understating} error mass. The aggravating detail is that
training increases both $\mu$ and $\sigma$, so the gap $e^{\sigma^2/2}$ is
largest precisely when the signal is weakest; the two effects compound.

\begin{proposition}[Clip-and-renormalize is not Lipschitz; projection is]
\label{prop:lip}
Let $R(y) = n\pos{y}/\|\pos{y}\|_1$ denote clip-and-renormalize, as used
in~\eqref{eq:rlu} and after~\eqref{eq:zlg}, and let $\Pi_{\simplex}$ denote
Euclidean projection onto $\simplex$. Then
\[
 \sup_{a \ne b} \frac{\|R(a)-R(b)\|}{\|a-b\|} = \infty,
 \qquad\text{whereas}\qquad
 \|\Pi_{\simplex}(a) - \Pi_{\simplex}(b)\| \le \|a-b\| \quad \forall a,b .
\]
\end{proposition}

\begin{proof}
Projection onto a nonempty closed convex set is non-expansive, which gives
the second claim. For the first, take $a = (\delta, \delta/2, -1, \ldots)$
and $b = a + \eta u$ with $u$ supported on the first two coordinates. As
$\delta \to 0$ the denominator $\|\pos{a}\|_1 \to 0$, so $R$ magnifies the
perturbation without bound.
\end{proof}

The practical consequence is that the final normalization step of RLU and
ZLG$^+$ can \emph{amplify} gradient noise, while the corresponding step of
VLI$^*$ (Section~\ref{sec:vli}) cannot. Section~\ref{sec:numerics} reports a
measured amplification factor exceeding $4\times 10^4$ for $R$ against
exactly $1$ for $\Pi_{\simplex}$.

\subsection{iLRG: a conditioning failure, not a calibration failure}

iLRG uses no auxiliary statistic, so Theorem~\ref{thm:collapse} does not
apply. Its difficulty is different. By Lemma~\ref{lem:identities},
$G_j = -a_j + E_j$, so as $\bar\epsilon \to 0$ both $a_j$ and $E_j$ vanish
and the bias gradient used as a \emph{denominator} in
$\hat e_j = \nabla_{W_j}/\nabla_{b_j}$ tends to zero. The ratio remains
finite in exact arithmetic but its relative error is amplified without
bound; the reference implementation must clamp the exponent of the ratio and
substitute a magic value when the denominator is exactly zero. The
subsequent linear system is also increasingly ill-conditioned, since
$p_i \to e_{y_i}$ makes all recovered probability vectors collapse towards
one-hot vectors.

\subsection{Summary of baseline diagnoses}

\begin{center}
\begin{tabular}{llp{6.3cm}}
\toprule
Method & Mechanism & Limit or symptom \\
\midrule
iLRG    & conditioning       & division by a vanishing bias gradient; near-singular pseudo-inverse \\
LLG     & none (scale-free)  & survives; but assumes $r_j \equiv 1$, uniform $Q$, and multiplies a single-batch answer by $m$ \\
LLG$^+$ & Thm~\ref{thm:collapse} + Prop~\ref{prop:gap} & auxiliary offset dominates $G$; worse than LLG \\
ZLG$^+$ & Thm~\ref{thm:collapse} + Prop~\ref{prop:gap} & $\widehat N \to n\bar p \approx$ uniform prior \\
RLU     & Thm~\ref{thm:collapse} + Props~\ref{prop:gap},~\ref{prop:jensen},~\ref{prop:lip} & $\widehat N \to n\,\mathrm{shift}$; Jensen bias $e^{\sigma^2/2}$; non-Lipschitz normalization \\
\bottomrule
\end{tabular}
\end{center}

% ============================================================================
\section{VLI and VLI$^*$}
\label{sec:vli}

\subsection{The design principle}

Lemma~\ref{lem:identities} writes $a_j = N_j r_j \bar\epsilon$, so the
label information lies in the \emph{relative} structure of $G$, while
$\bar\epsilon$ --- the quantity that vanishes under training --- enters only
as a common factor. This suggests requiring of an estimator exactly the
property the baselines lack.

\begin{proposition}[Scale invariance]
\label{prop:invariance}
Let $\Psi$ denote the map defined in Theorem~\ref{thm:closed}, and let
$\Psi_{\mathrm{gen}}$ denote the map of Theorem~\ref{thm:general}. For every
$c>0$,
\[
 \Psi(cG) = \Psi(G), \qquad \Psi_{\mathrm{gen}}(cG, r, Q) = \Psi_{\mathrm{gen}}(G,r,Q) .
\]
Consequently neither estimator can exhibit the collapse of
Theorem~\ref{thm:collapse}: no quantity of absolute magnitude enters either
map.
\end{proposition}

\begin{proof}
In Theorem~\ref{thm:closed} the signal appears only through the ratio
$G_j/G_{\max}$, and $G_{\max}(cG) = c\,G_{\max}(G)$ for $c>0$. In
Theorem~\ref{thm:general} the solution $a$ of~\eqref{eq:system} is
positively homogeneous in $G$, and the output is the ratio
$n(a_j/r_j)/\sum_k (a_k/r_k)$.
\end{proof}

\begin{remark}[The heuristic version is only partly invariant]
The original VLI procedure builds a zero-sum basis with diagonal
$d = n^{-1}\sum_{j:G_j<0}G_j$ and off-diagonal $-d/(C-1)$, scales $d$ by a
``dynamic boost'' $f_1 f_2 f_3$, and greedily peels basis columns from a
residual initialized at $G$. The factor $f_1$ is a function of normalized
projections and is scale-invariant, but $f_2 = (w_2\|G\|_2+1)^{-1}$ and
$f_3 = w_3^{-100|d|}$ depend on the absolute magnitude of $G$. The heuristic
therefore violates Proposition~\ref{prop:invariance}, and its five constants
must be retuned per dataset, model and accuracy level.
Theorem~\ref{thm:closed} removes them.
\end{remark}

\subsection{The symmetry assumption, and why it is forced}

\begin{assumption}[Symmetric error geometry]
\label{ass:sym}
All present classes have equal mean error, $r_j = 1$, and every class
distributes its error mass uniformly over the other classes,
$Q_{jk} = 1/(C-1)$ for $k \ne j$. At least one class is absent from the
window, and $\bar\epsilon > 0$.
\end{assumption}

Assumption~\ref{ass:sym} is false for any real model. The following
proposition explains why it is nonetheless the correct choice for an
estimator that is required to use no class-specific side information.

\begin{proposition}[Uniqueness under label exchangeability]
\label{prop:exchange}
The uniform matrix $Q_{jk} = 1/(C-1)$ $(k\ne j)$, $Q_{jj}=0$, is the
\emph{unique} row-stochastic matrix with zero diagonal satisfying
$P Q P^T = Q$ for every $C\times C$ permutation matrix $P$.
\end{proposition}

\begin{proof}
Permutation invariance means $Q_{\sigma(j)\sigma(k)} = Q_{jk}$ for every
permutation $\sigma$. Given any two off-diagonal positions $(j,k)$ and
$(j',k')$ with $j\ne k$, $j'\ne k'$, there is a permutation carrying one to
the other, so all off-diagonal entries are equal, say to $q$. Row
stochasticity with $C-1$ off-diagonal entries per row forces
$q = 1/(C-1)$. Conversely the uniform matrix is clearly invariant.
\end{proof}

An attacker who does not know which classes are confusable has no basis for
breaking the symmetry among class labels; any non-uniform $Q$ asserts
class-specific knowledge. Proposition~\ref{prop:exchange} says the uniform
choice is then not a simplification but the only admissible one --- the same
style of argument that derives a uniform prior from exchangeability. A
degrees-of-freedom count agrees: system~\eqref{eq:system} provides $C$
equations, while the unknowns are $a$ together with the $C(C-2)$ free
parameters of $Q$; identifiability with no side information requires removing
all of the latter, and the uniform matrix is the unique zero-parameter,
label-equivariant choice that keeps the system solvable in closed form.

\subsection{The closed form}

\begin{theorem}[Closed-form VLI$^*$]
\label{thm:closed}
Under Assumption~\ref{ass:sym},
\begin{equation}
 G_j = \frac{n\bar\epsilon - C a_j}{C-1},
 \qquad
 G_{\max} = \frac{n\bar\epsilon}{C-1},
 \qquad
 \boxed{\;N_j = \frac{n}{C}\Bigl(1 - \frac{G_j}{G_{\max}}\Bigr)\;}
 \label{eq:closed}
\end{equation}
and the final expression is invariant under $G \mapsto cG$ for $c>0$.
\end{theorem}

\begin{proof}
Under uniform confusion
$(Q^Ta)_j = \bigl(\sum_k a_k - a_j\bigr)/(C-1)$, so
by~\eqref{eq:system}
\[
 G_j = -a_j + \frac{\sum_k a_k - a_j}{C-1}
     = \frac{\sum_k a_k - C a_j}{C-1} ,
\]
and $\sum_k a_k = n\bar\epsilon$ gives the first equality. Since
$a_j \ge 0$ and $G_j$ is affine and strictly decreasing in $a_j$ with the
same slope $-C/(C-1)$ for every $j$, the maximum of $G$ is attained exactly
where $a_j = 0$, that is at an absent class, giving
$G_{\max} = n\bar\epsilon/(C-1)$; Assumption~\ref{ass:sym} guarantees such a
class exists. Dividing,
\[
 \frac{G_j}{G_{\max}} = 1 - \frac{C a_j}{n\bar\epsilon}
 \quad\Longrightarrow\quad
 a_j = \frac{n\bar\epsilon}{C}\Bigl(1 - \frac{G_j}{G_{\max}}\Bigr),
\]
and $r_j = 1$ gives $a_j = N_j\bar\epsilon$, so $\bar\epsilon$ cancels and
\eqref{eq:closed} follows. A common positive factor cancels in
$G_j/G_{\max}$.
\end{proof}

The cancellation of $\bar\epsilon$ in the last step is the entire point. The
role of $G_{\max}$ is to serve as an \emph{in-window thermometer} for
$\bar\epsilon$: an absent class has $a_j = 0$, so its coordinate reads the
error scale directly off the attacked update. Where RLU, ZLG$^+$ and LLG$^+$
borrow that scale from auxiliary data and are punished by
Proposition~\ref{prop:gap}, VLI$^*$ reads it from inside the window.

\begin{proposition}[Structural properties of~\eqref{eq:closed}]
\label{prop:props}
Write $\mathrm{raw}_j = \tfrac{n}{C}(1 - G_j/G_{\max})$ with $G_{\max}>0$.
Then
\begin{enumerate}
\item $\mathrm{raw}_j \ge 0$ for every $j$, with equality at the maximizing
  coordinate; no clipping is required.
\item $\sum_j \mathrm{raw}_j = n$, by~\eqref{eq:zerosum}; no renormalization
  is required.
\item Under Assumption~\ref{ass:sym},
  $\mathrm{raw}_j - \mathrm{raw}_k = N_j - N_k$ exactly.
\end{enumerate}
\end{proposition}

\begin{proof}
(i) $G_j \le G_{\max}$ and $G_{\max}>0$ give $G_j/G_{\max} \le 1$.
(ii) $\sum_j \mathrm{raw}_j = \tfrac{n}{C}\bigl(C - G_{\max}^{-1}\sum_j G_j\bigr) = n$.
(iii) From Theorem~\ref{thm:closed},
$G_k - G_j = C\bar\epsilon(N_j-N_k)/(C-1)$ and
$G_{\max} = n\bar\epsilon/(C-1)$, hence
$\mathrm{raw}_j - \mathrm{raw}_k = \tfrac{n}{C}\cdot\tfrac{G_k-G_j}{G_{\max}}
 = \tfrac{n}{C}\cdot\tfrac{C(N_j-N_k)}{n} = N_j - N_k$.
\end{proof}

Items (i) and (ii) remove exactly the two steps that
Proposition~\ref{prop:lip} identifies as pathological in RLU and ZLG$^+$.
Item (iii) determines the correct way to repair a perturbed estimate: the
label information lies in the \emph{differences} of $\mathrm{raw}$, so the
repair must be \emph{additive}. Euclidean projection onto $\simplex$ has
exactly this form,
\[
 \Pi_{\simplex}(y)_j = \max(y_j - \theta, 0),
 \qquad \theta \text{ chosen so that } \textstyle\sum_j \max(y_j-\theta,0) = n,
\]
obtained from the Karush--Kuhn--Tucker conditions of
$\min_{x\in\simplex}\tfrac12\|x-y\|^2$: it subtracts a single threshold and
clips, preserving differences among surviving coordinates. By
Proposition~\ref{prop:lip} it is also non-expansive. A multiplicative
rescaling preserves ratios instead and is therefore the wrong repair for
counts.

\subsection{A sufficient condition weaker than Assumption~\ref{ass:sym}}

\begin{proposition}[Weaker sufficient condition]
\label{prop:weak}
Suppose $r_j = 1$ on present classes, at least one class is absent, and the
incoming error mass satisfies
\begin{equation}
 (Q^Ta)_j = \frac{\sum_k a_k - a_j}{C-1} \qquad \text{for every } j .
 \label{eq:flow}
\end{equation}
Then~\eqref{eq:closed} holds exactly. Uniform $Q$ implies~\eqref{eq:flow},
but~\eqref{eq:flow} is strictly weaker: it constrains the vector $Q^Ta$
rather than the individual entries of $Q$.
\end{proposition}

\begin{proof}
Substituting~\eqref{eq:flow} into $G = -a + Q^Ta$ reproduces the expression
for $G_j$ in Theorem~\ref{thm:closed}, from which the rest of that proof
applies verbatim. Strictness follows from the example below.
\end{proof}

Condition~\eqref{eq:flow} is a statement about the \emph{aggregate} inflow
into each class and can hold for highly non-uniform $Q$. Concretely, take
$C=10$ with present classes carrying equal counts and let $Q$ consist of
confusable pairs $(0,1), (2,3), \ldots$, each class sending a fraction $s$ of
its error mass to its partner and spreading the rest uniformly. If the absent
classes form a complete pair, then $G$ takes only two distinct values --- one
across all present classes and one across all absent classes --- so $G$
remains affine in $N$ and recovery is exact \emph{for every} $s$, including
$s = 0.9$. Section~\ref{sec:numerics} verifies this. The practical reading is
that the exactness region of~\eqref{eq:closed} is a set containing the
uniform-$Q$ point, not that single point, and that departures are governed by
the residual of~\eqref{eq:flow} rather than by the distance of $Q$ from
uniform.

\subsection{Known geometry: the oracle upper bound}

\begin{assumption}[Known irreducible error geometry]
\label{ass:geom}
$Q$ is known, row-stochastic, nonnegative, zero-diagonal and irreducible;
$r$ is known and strictly positive on present classes; at least one class is
absent; $\bar\epsilon>0$.
\end{assumption}

\begin{theorem}[Recovery with known $(r,Q)$]
\label{thm:general}
Under Assumption~\ref{ass:geom}, let $v>0$ be the normalized Perron vector
with $Q^Tv=v$, $\one^Tv=1$. Every solution of~\eqref{eq:system} is
$a = a_0 + tv$ with $a_0 = (I-Q^T)^{+}(-G)$, and the requirements $a \ge 0$
together with the existence of an absent class select
\begin{equation}
 t = \max_j\Bigl(-\frac{a_{0,j}}{v_j}\Bigr),
 \qquad
 \widehat N_j = n\,\frac{a_j/r_j}{\sum_k a_k/r_k}.
 \label{eq:general}
\end{equation}
For exact $G$, $r$ and $Q$, \eqref{eq:general} equals the true counts.
\end{theorem}

\begin{proof}
$Q$ irreducible and stochastic gives, by Perron--Frobenius, a
one-dimensional nullspace of $I-Q^T$ spanned by a strictly positive $v$, so
all solutions differ by $tv$. The constraints $a_{0,j}+tv_j\ge0$ are
equivalent to $t \ge \max_j(-a_{0,j}/v_j)$, and at the smallest feasible $t$
at least one coordinate vanishes. The true $a$ is nonnegative and has a zero
coordinate because some class is absent, so it is this boundary solution. For
present classes $a_j = N_jr_j\bar\epsilon$, hence
$a_j/r_j = N_j\bar\epsilon$, and normalizing these to sum to $n$ returns the
counts. Absent classes contribute zero regardless of the positive placeholder
used for their $r_j$.
\end{proof}

\begin{corollary}[Difficulty-mismatch bound]
\label{cor:kappa}
Suppose $Q$ is known and the boundary solution $a$ is exact, but the
estimator substitutes $r_j \equiv 1$. With
$\kappa = r_{\max}/r_{\min}$ over present classes, the instance recovery
rate $\sum_j \min(\widehat\pi_j, \pi_j)$ is at least $1/\kappa$.
\end{corollary}

\begin{proof}
The estimate is $q_j = a_j/A$ with $A=\sum_k a_k$, while
$\pi_j = (a_j/r_j)/Z$ with $Z = \sum_k a_k/r_k$. Hence
$q_j/\pi_j = r_j Z/A$, and $Z/A$ is a weighted average of the $1/r_k$, so
$Z/A \ge 1/r_{\max}$. Thus $q_j \ge \pi_j/\kappa$, so
$\min(q_j,\pi_j) \ge \pi_j/\kappa$; summing gives the claim.
\end{proof}

Theorem~\ref{thm:general} is \emph{not} an attack: $r$ and $Q$ are oracle
inputs. Its role is to supply an upper bound against which
Theorem~\ref{thm:closed} can be measured, with
Corollary~\ref{cor:kappa} bounding the gap. Together the three results form a
ladder: the closed form is the deployable estimator, the generalized form is
the ceiling, and the corollary certifies that the distance between them
degrades gracefully rather than collapsing.

% ============================================================================
\section{How VLI$^*$ clears each barrier}

\begin{center}
\begin{tabular}{p{4.3cm}p{4.0cm}p{5.2cm}}
\toprule
Barrier & Who it stops & How VLI$^*$ avoids it \\
\midrule
Lemma~\ref{lem:norm}: $\|G\|_1 \le 2nL_{\mathrm{train}}$
  & every method
  & only the ratio $G_j/G_{\max}$ is used (Prop.~\ref{prop:invariance}); the
    common factor cancels \\
\addlinespace
Thm~\ref{thm:collapse}: $\widehat\pi \to \Phi(c)$
  & LLG$^+$, ZLG$^+$, RLU
  & no external $c$ exists; the error scale is read from $G_{\max}$ inside
    the window \\
\addlinespace
Prop.~\ref{prop:gap}: $\bar\epsilon_{\mathrm{aux}}/\bar\epsilon_{\mathrm{train}}\to\infty$
  & LLG$^+$, ZLG$^+$, RLU
  & no auxiliary magnitude is referenced; the optional $(r,Q)$ of
    Thm~\ref{thm:general} are normalized, so $\bar\epsilon$ cancels from them
    by construction \\
\addlinespace
Prop.~\ref{prop:jensen}: Jensen gap $e^{\sigma^2/2}$
  & RLU
  & no softmax of averaged logits is formed \\
\addlinespace
Prop.~\ref{prop:lip}: $R$ not Lipschitz
  & RLU, ZLG$^+$
  & clipping and renormalization are unnecessary
    (Prop.~\ref{prop:props}(i)--(ii)); the defensive repair
    $\Pi_{\simplex}$ is non-expansive and additive, matching
    Prop.~\ref{prop:props}(iii) \\
\addlinespace
Vanishing denominator $\nabla W/\nabla b$
  & iLRG
  & the bias gradient is never used as a divisor \\
\addlinespace
Hand-tuned constants
  & VLI heuristic
  & the signal-matched diagonal
    $\mathrm{boost}^\star = (C-1)G_{\max}/\sum_j\pos{-G_j}$ is derived, and
    under Assumption~\ref{ass:sym} peeling with it reproduces
    \eqref{eq:closed} \\
\bottomrule
\end{tabular}
\end{center}

The trade being made should be stated plainly. The baselines model the error
geometry \emph{better} --- they measure confusion and posteriors from data ---
but pay for it with scale dependence. VLI$^*$ models the geometry
\emph{worse}, by Assumption~\ref{ass:sym}, but is scale-free. The claim of
this report is not that Assumption~\ref{ass:sym} is true. It is that

\begin{quote}
a misspecified but scale-invariant model outperforms a better-specified but
scale-dependent one, because geometry misspecification is \emph{bounded}
(Corollary~\ref{cor:kappa}, Proposition~\ref{prop:weak}) whereas scale
collapse is \emph{total} (Theorem~\ref{thm:collapse}).
\end{quote}

This is structurally the argument that explains Naive Bayes: its conditional
independence assumption is essentially always violated, yet the classifier
performs well because the decision rule depends only on an ordering, not on
calibrated probabilities. Here, by Proposition~\ref{prop:props}(iii), the
estimator depends only on differences of $\mathrm{raw}$, so a wrong $Q$ is
harmless as long as the affine relation between $G$ and $N$ is preserved.

% ============================================================================
\section{Numerical verification}
\label{sec:numerics}

\paragraph{What these tables are.} Every number below is computed on the
generative model implied by Lemma~\ref{lem:identities}: given counts $N$, a
difficulty vector $r$, a confusion matrix $Q$ and a scale $\bar\epsilon$, we
form $a_j = N_jr_j\bar\epsilon$, set $G = -(I-Q^T)a$ exactly, and pass $G$ to
the implementation. They therefore verify the algebra of
Theorems~\ref{thm:closed} and~\ref{thm:general} and measure sensitivity to
controlled violations of Assumption~\ref{ass:sym}. \textbf{They are not
experiments on a trained network} and must not be read as such; an empirical
study requires checkpoints at several accuracy levels and is future work
(Section~\ref{sec:limits}).

\paragraph{Exactness and scale invariance.}
With $C=4$, $N=(6,4,2,0)$, uniform $Q$, $r\equiv1$, the recovered counts are
$(6,4,2,0)$ exactly, and remain so as $\bar\epsilon$ --- hence the whole
signal --- is driven to zero over seven orders of magnitude, confirming
Proposition~\ref{prop:invariance}.

\begin{center}
\begin{tabular}{lll}
\toprule
$\bar\epsilon$ & $\|G\|_1$ & recovered counts \\
\midrule
$5\times10^{-1}$ & $5.33\times10^{0}$  & $(6,4,2,0)$ \\
$1\times10^{-1}$ & $1.07\times10^{0}$  & $(6,4,2,0)$ \\
$1\times10^{-3}$ & $1.07\times10^{-2}$ & $(6,4,2,0)$ \\
$1\times10^{-8}$ & $1.07\times10^{-7}$ & $(6,4,2,0)$ \\
\bottomrule
\end{tabular}
\end{center}

\paragraph{Proposition~\ref{prop:lip} measured.} Maximum observed
amplification $\|f(a)-f(b)\|/\|a-b\|$ over $2\times10^5$ random pairs:

\begin{center}
\begin{tabular}{lr}
\toprule
Map & Maximum amplification \\
\midrule
$\Pi_{\simplex}$ (projection, used by VLI$^*$) & $1.00$ \\
$R$ (clip-and-renormalize, used by RLU, ZLG$^+$) & $44\,437.19$ \\
\bottomrule
\end{tabular}
\end{center}

\paragraph{Structured confusion.} $C=10$,
$N=(120,90,70,50,30,20,10,5,0,0)$, $r\equiv1$, and $Q$ built from confusable
pairs with strength $s$ ($s=1/(C-1)\approx0.111$ is uniform). This is the
most consequential violation, since real models do confuse specific pairs.

\begin{center}
\begin{tabular}{lr}
\toprule
pair strength $s$ & instance recovery \\
\midrule
$0.111$ (uniform) & $1.000$ \\
$0.200$ & $0.982$ \\
$0.300$ & $0.962$ \\
$0.500$ & $0.884$ \\
$0.700$ & $0.734$ \\
$0.900$ & $0.638$ \\
\bottomrule
\end{tabular}
\end{center}

\paragraph{Proposition~\ref{prop:weak} verified.} With the same $C=10$ and
pair-structured $Q$, but present classes carrying equal counts and the two
absent classes forming a complete pair, recovery is exact for every strength
tested, despite $Q$ being maximally non-uniform:

\begin{center}
\begin{tabular}{lrr}
\toprule
configuration & $s=0.3$ & $s=0.9$ \\
\midrule
balanced present, absent pair intact & $1.000$ & $1.000$ \\
balanced present, pairs split across present/absent & $0.964$ & $0.864$ \\
mildly imbalanced present, pair intact & $0.985$ & $0.667$ \\
strongly imbalanced present, pair intact & $0.940$ & $0.530$ \\
\bottomrule
\end{tabular}
\end{center}

\paragraph{The residual of~\eqref{eq:flow} predicts the loss.} Over $400$
random configurations (varying $N$, number of present classes, and
$Q \in \{\text{uniform}, \text{paired}, \text{cyclic}\}$), let
$\mathrm{defect} = \|Q^Ta - (\one^Ta\,\one - a)/(C-1)\|_2/\|a\|_2$. The
correlation between defect and recovery is $-0.873$, and a zero defect was
exact in all $138$ cases where it occurred, as Proposition~\ref{prop:weak}
predicts.

\begin{center}
\begin{tabular}{lrr}
\toprule
defect range & count & mean instance recovery \\
\midrule
$0$ (exactly) & $138$ & $1.000$ \\
$(0,\,0.05)$  & $17$  & $0.958$ \\
$[0.05,\,0.1)$ & $39$ & $0.917$ \\
$[0.1,\,0.3)$ & $65$  & $0.832$ \\
$[0.3,\,\infty)$ & $141$ & $0.716$ \\
\bottomrule
\end{tabular}
\end{center}

\paragraph{Difficulty spread, against Corollary~\ref{cor:kappa}.} Uniform
$Q$, $r$ linearly spread, estimator using $r\equiv1$. The corollary's bound
is correct but very loose; the realized degradation is far gentler.

\begin{center}
\begin{tabular}{lrr}
\toprule
$\kappa$ & instance recovery & bound $1/\kappa$ \\
\midrule
$1.00$  & $1.000$ & $1.000$ \\
$1.39$  & $0.965$ & $0.720$ \\
$1.78$  & $0.934$ & $0.562$ \\
$3.33$  & $0.853$ & $0.300$ \\
$8.00$  & $0.744$ & $0.125$ \\
$39.11$ & $0.635$ & $0.026$ \\
\bottomrule
\end{tabular}
\end{center}

\paragraph{Perturbation of the anchor.} Gaussian noise of standard deviation
$\sigma\cdot\|G\|_\infty$ added to $G$, $2000$ repetitions. The last column
reports how often the noise moves $\argmax_j G_j$ onto a \emph{present}
class, which destroys the thermometer interpretation of $G_{\max}$.

\begin{center}
\begin{tabular}{lrr}
\toprule
$\sigma$ & instance recovery & anchor lands on a present class \\
\midrule
$0\%$  & $1.000$ & $0.0\%$ \\
$1\%$  & $0.991$ & $0.0\%$ \\
$5\%$  & $0.954$ & $8.6\%$ \\
$10\%$ & $0.914$ & $22.6\%$ \\
$25\%$ & $0.820$ & $44.0\%$ \\
$50\%$ & $0.722$ & $53.4\%$ \\
\bottomrule
\end{tabular}
\end{center}

\paragraph{Absent-class requirement.} When no class is absent the formula is
forced to zero its $\argmax$ coordinate, so the smallest present class is
sacrificed. The penalty therefore scales with that class's share of $n$, not
with absence itself: for the long-tailed $C=10$ configuration above, removing
the absent classes entirely still yields $0.975$, whereas a small
near-balanced window suffers considerably more. A perfectly balanced window
under Assumption~\ref{ass:sym} gives $G \equiv 0$, and the implementation
correctly refuses rather than returning a fabricated answer.

% ============================================================================
\section{Limitations and threats to validity}
\label{sec:limits}

\begin{enumerate}
\item \textbf{$(r,Q)$ are oracle inputs.} Theorem~\ref{thm:general} assumes
  them known. That $\bar\epsilon$ cancels from both by construction
  (Section~\ref{sec:notation}) is a reason to expect them to transfer from
  auxiliary to training data better than the absolute magnitudes used by the
  baselines, but this is a motivated hypothesis, not a theorem, and must be
  measured. The gap between $0.638$ and $1.000$ in the
  structured-confusion table is the quantitative value of solving it.

\item \textbf{$\kappa$ grows with training.} Since $r_j = \bar\epsilon_j/\bar\epsilon$
  is a ratio, $\bar\epsilon$ cancels, but the \emph{spread} of $r$ need not
  be stable: easy classes may reach $\bar\epsilon_j\approx10^{-6}$ while hard
  ones remain at $10^{-2}$, giving $\kappa \approx 10^4$ where an
  undertrained model has $\kappa \approx 2$. The regime in which VLI$^*$
  wins is thus also the regime in which Assumption~\ref{ass:sym} is least
  accurate. This is the central tension of the approach and should be
  reported, not hidden; it is directly measurable from checkpoints at
  increasing accuracy.

\item \textbf{The anchor is a single coordinate.} $G_{\max}$ carries the
  whole scale calibration, so moderate additive noise --- of the kind a
  differentially private or compressed update would introduce --- degrades
  VLI$^*$ more than it would an estimator averaging over coordinates. This
  suggests the natural defence: perturbation sufficient to hide which class
  is absent, rather than strong privacy noise.

\item \textbf{Reading the gradient is a white-box convenience.} The
  identities require the cross-entropy gradient. An implementation that reads
  the parameter's gradient buffer before the optimizer step obtains exactly
  that, but a real server observes an accumulated weight update, which under
  momentum, weight decay, or several local steps is not $\sum_i(p_i-e_{y_i})$.
  Zero momentum is a requirement of the derivation, not a convenience.

\item \textbf{Proximal and dynamic regularizers break the signal.} A FedProx
  term contributes $\mu(b-b_g)$ to the bias gradient, which is neither
  zero-sum nor of the form~\eqref{eq:system}; both~\eqref{eq:zerosum} and the
  $G_{\max}$ anchor then fail. The closed form should be restricted to plain
  FedAvg, and the heuristic's apparent operation under FedProx has no
  theoretical basis.

\item \textbf{Aggregation.} If the server observes only an average over
  several clients, all methods in this report recover the union of their
  labels, not a per-client decomposition.

\item \textbf{Numerics are model-based.} As stated in
  Section~\ref{sec:numerics}, the tables evaluate the generative model, not a
  trained network. The empirical programme that would complete this report
  is: produce checkpoints spanning a range of training accuracies; measure
  $\bar\epsilon_{\mathrm{train}}$, $\bar\epsilon_{\mathrm{aux}}$, $\kappa$
  and $Q$ on each; verify Lemma~\ref{lem:norm} and
  Proposition~\ref{prop:gap} directly; and plot instance recovery against
  $\bar\epsilon_{\mathrm{train}}$ for all methods, reporting separately the
  fraction of windows containing an absent class.
\end{enumerate}

% ============================================================================
\begin{thebibliography}{9}
\bibitem{ilrg} J.~Ma et al.
  \emph{Instance-wise Batch Label Restoration via Gradients in Federated
  Learning}. ICLR, 2023.
\bibitem{llg} A.~Wainakh et al.
  \emph{User-Level Label Leakage from Gradients in Federated Learning}.
  arXiv:2105.09369, 2021.
\bibitem{zlg} J.~Geng et al.
  \emph{Towards General Deep Leakage in Federated Learning}.
  arXiv:2110.09074, 2021.
\bibitem{rlu} H.~Chen and H.~Vikalo.
  \emph{Recovering Labels from Local Updates in Federated Learning}.
  ICML, 2024.
\bibitem{duchi} J.~Duchi, S.~Shalev-Shwartz, Y.~Singer and T.~Chandra.
  \emph{Efficient Projections onto the $\ell_1$-Ball for Learning in High
  Dimensions}. ICML, 2008.
\bibitem{nb} P.~Domingos and M.~Pazzani.
  \emph{On the Optimality of the Simple Bayesian Classifier under Zero-One
  Loss}. Machine Learning, 29:103--130, 1997.
\end{thebibliography}

\end{document}
